% !TeX TS-program = pdflatex


\documentclass[a4paper]{article}

% \usepackage[default]{fontsetup}

\usepackage{fancyhdr}
\usepackage{extramarks}
\usepackage{amsmath}
\usepackage{amsthm}
\usepackage{amsfonts}
\usepackage{tikz}
\usepackage[plain]{algorithm}
\usepackage{algpseudocode}
\usepackage{enumerate}
\usepackage{tikz}

\usetikzlibrary{automata,positioning}

%
% Basic Document Settings
%  

\topmargin=-0.2in
\evensidemargin=0in
\oddsidemargin=0in
\textwidth=6.5in
\textheight=9.5in
\headsep=0.25in

\linespread{1.1}

\pagestyle{fancy}
\lhead{\hmwkAuthorName}
\chead{\hmwkClass : \hmwkTitle}
\rhead{\firstxmark}
\lfoot{\lastxmark}
\cfoot{\thepage}

\renewcommand\headrulewidth{0.4pt}
\renewcommand\footrulewidth{0.4pt}

\setlength\parindent{0pt}

%
% Create Problem Sections
%

\newcommand{\enterProblemHeader}[1]{
    \nobreak\extramarks{}{Problem \arabic{#1} continued on next page\ldots}\nobreak{}
    \nobreak\extramarks{Problem \arabic{#1} (continued)}{Problem \arabic{#1} continued on next page\ldots}\nobreak{}
}

\newcommand{\exitProblemHeader}[1]{
    \nobreak\extramarks{Problem \arabic{#1} (continued)}{Problem \arabic{#1} continued on next page\ldots}\nobreak{}
    \stepcounter{#1}
    \nobreak\extramarks{Problem \arabic{#1}}{}\nobreak{}
}

\newcommand*\circled[1]{\tikz[baseline=(char.base)]{
		\node[shape=circle,draw,inner sep=2pt] (char) {#1};}}


\setcounter{secnumdepth}{0}
\newcounter{partCounter}
\newcounter{homeworkProblemCounter}
\setcounter{homeworkProblemCounter}{1}
\nobreak\extramarks{Problem \arabic{homeworkProblemCounter}}{}\nobreak{}

%
% Homework Problem Environment
%
% This environment takes an optional argument. When given, it will adjust the
% problem counter. This is useful for when the problems given for your
% assignment aren't sequential. See the last 3 problems of this template for an
% example.
%

\newenvironment{homeworkProblem}[1][-1]{
    \ifnum#1>0
        \setcounter{homeworkProblemCounter}{#1}
    \fi
    \section{Problem \arabic{homeworkProblemCounter}}
    \setcounter{partCounter}{1}
    \enterProblemHeader{homeworkProblemCounter}
}{
    \exitProblemHeader{homeworkProblemCounter}
}

%
% Homework Details
%   - Title
%   - Class
%   - Due date
%   - Name
%   - Student ID

\newcommand{\hmwkTitle}{Homework\ \#03}
\newcommand{\hmwkClass}{Probability \& Statistics for EECS}
\newcommand{\hmwkDueDate}{To be announced}
\newcommand{\hmwkAuthorName}{Jianxuan Yao}
\newcommand{\hmwkAuthorID}{2025521024}


%
% Title Page
%

\title{
    \vspace{2in}
    \textmd{\textbf{\hmwkClass:\\  \hmwkTitle}}\\
    \normalsize\vspace{0.1in}\small{Due\ on\ \hmwkDueDate\ at 23:59}\\
	\vspace{4in}
}

\author{
	Name: \textbf{\hmwkAuthorName} \\
	Student ID: \hmwkAuthorID}
\date{}

\renewcommand{\part}[1]{\textbf{\large Part \Alph{partCounter}}\stepcounter{partCounter}\\}

%
% Various Helper Commands
%

% Useful for algorithms
\newcommand{\alg}[1]{\textsc{\bfseries \footnotesize #1}}
% For derivatives
\newcommand{\deriv}[1]{\frac{\mathrm{d}}{\mathrm{d}x} (#1)}
% For partial derivatives
\newcommand{\pderiv}[2]{\frac{\partial}{\partial #1} (#2)}
% Integral dx
\newcommand{\dx}{\mathrm{d}x}
% Alias for the Solution section header
\newcommand{\solution}{\textbf{\large Solution}}
% Probability commands: Expectation, Variance, Covariance, Bias
\newcommand{\E}{\mathrm{E}}
\newcommand{\Var}{\mathrm{Var}}
\newcommand{\Cov}{\mathrm{Cov}}
\newcommand{\Bias}{\mathrm{Bias}}

\newcommand{\Prob}{P}

\begin{document}


% \maketitle
% \thispagestyle{empty}
% \pagebreak

\date{
Due date: Due on Oct. 11, 2026, 23:59 UTC+8}
\title{SI 140A.02  Probability \& Statistics for EECS, Fall 2026 \\
Homework 3}
\maketitle
Read all the instructions below carefully before you start working on the assignment, and before you make a submission.
\begin{itemize}
    \item You are required to write down all the major steps towards making your conclusions; otherwise you may obtain limited points of the problem.
    \item Write your homework in English; otherwise you will get no points of this homework.
    \item Any form of plagiarism will lead to $0$ point of this homework. 
    \item You are required to submit your homework in PDF format generated by LaTeX, otherwise you will get no points of this homework.
\end{itemize}
\newpage

\begin{homeworkProblem}[1]
% Source: Fall 2025 HW3, Problem 2(b), PDF page 3.
There are two states, Red and Blue, initially with 100 residents each. Every resident supports exactly one of two parties, Democrat or Republican. Initially, 60 residents of Blue and 30 residents of Red are Democrats. Ten residents move from Blue to Red, and nobody changes parties. Let $k$ be the number of Democrats among the ten people who move.

Whenever a resident is selected, the selection is uniform over all 200 people. Let $D$ be the event that the selected resident is a Democrat, and let $B$ be the event that the resident lives in Blue at the time of selection.
\begin{enumerate}[(a)]
\item Find $\Prob(D\mid B)$, $\Prob(D\mid B^c)$, and $\Prob(D)$ before the move. Does the move change $\Prob(D)$? Explain.
\item Determine all integer values of $k$ for which the proportion of Democrats increases strictly in \emph{both} states. Explain, using the law of total probability with the appropriate weights, how both proportions can increase while the overall proportion stays unchanged.
\item For this part, take $k=5$ and consider the population after the move. A survey interviews the selected resident only if their name appears on a list containing all Blue residents and all Democrats in Red. Let $A$ be the event that the selected resident is on this list. Find $\Prob(D\mid A)$ using the law of total probability conditional on $A$, with the partition $\{B,B^c\}$, and state the two group weights. Explain why this probability differs from the proportion of Democrats in the whole population.
\end{enumerate}

\textbf{Solution: }

\begin{enumerate}
    \item[(a)]
    
    Before the move:

    $$
    P(D \mid B) = 0.6, \quad P(D \mid B^c) = 0.3, \quad P(D) = 0.45
    $$

    After the move, the total number of Democrats remains the same, so $P(D)$ does not change.
    
    \item[(b)]
    
    After the move:

    $$
    P(D \mid B) = \frac{60 - k}{90}, \quad P(D \mid B^c) = \frac{30 + k}{110}
    $$

    $$
    \frac{60 - k}{90} > 0.6 \implies k < 6
    $$

    $$
    \frac{30 + k}{110} > 0.3 \implies k > 3
    $$

    Therefore, $k \in \left\{ 4, 5\right\}$.
    
    $$
    P(D) = P(D \mid B) P(B) + P(D \mid B^c) P(B^c)
    $$

    $$
    P_{\text{before}}(B) = \frac{100}{200} = \frac{1}{2}, \quad P_{\text{before}}(B^c) = \frac{100}{200} = \frac{1}{2}
    $$

    $$
    P(D) = 0.6 \times \frac{1}{2} + 0.3 \times \frac{1}{2} = 0.45
    $$

    However, after the move, though the proportions change, the weights also change.

    $$
    P_{\text{after}}(B) = \frac{90}{200} = 0.45, \quad P_{\text{after}}(B^c) = \frac{110}{200} = 0.55
    $$

    Because the proportion of Democratic Party members in blue states (about 0.6) is significantly higher than that in red states (about 0.3), the flow of population from blue states to red states is equivalent to shifting the weight from the high proportion of Democratic Party members to the low proportion of Democratic Party members.

    \item[(c)]
    
    After the move, let $k=5$:

    $$
    P(D \mid B) = \frac{60 - 5}{90} = \frac{11}{18}, \quad P(D \mid B^c) = \frac{30 + 5}{110} = \frac{7}{22}
    $$

    $$
    P(B) = \frac{9}{20}, \quad P(B^c) = \frac{11}{20}
    $$

    $$
    P(D \mid A) = P(D \mid B, A) P(B \mid A) + P(D \mid B^c, A) P(B^c \mid A)
    $$

    $$
    P(B \mid A) = \frac{\vert B \cap A \vert}{\vert A \vert} = \frac{90}{90 + 35} = \frac{18}{25}, \quad P(B^c \mid A) = \frac{\vert B^c \cap A \vert}{\vert A \vert} = \frac{35}{90 + 35} = \frac{7}{25}
    $$

    $$
    P(D \mid B, A) = P(D \mid B) = \frac{11}{18}, \quad P(D \mid B^c, A) = 1
    $$

    $$
    \begin{aligned}
    \mathrm{Prob}(D \mid A)
    &= \left(\frac{11}{18}\right) \cdot \left(\frac{18}{25}\right) + (1) \cdot \left(\frac{7}{25}\right) \\
    &= \frac{11}{25} + \frac{7}{25} \\
    &= \frac{18}{25}
    \end{aligned}
    $$

    List $A$ contains all Democrats in Red, but it excludes the Republicans in Red. Therefore, the proportion of Democrats in list $A$ is higher than that in the whole population. 
\end{enumerate}
\end{homeworkProblem}

\newpage
\begin{homeworkProblem}[2]
% Source: Fall 2025 HW3, Problem 3, PDF page 4.
A sequence of $n\ge1$ mutually independent trials is performed. Each trial results in either success or failure, and the probability of success on trial $i$ is $p_i$, where $0\le p_i\le1$. Let $A_n$ be the event that the number of successes among the first $n$ trials is even.
\begin{enumerate}[(a)]
\item Find $\Prob(A_2)$ by listing the possible outcomes, and show that
\[
\Prob(A_2)=\frac{1+(1-2p_1)(1-2p_2)}{2}.
\]
\item By considering whether the last trial succeeds or fails, obtain a recursion for $\Prob(A_n)$ for $n\ge2$. Use the recursion to prove that
\[
\Prob(A_n)=\frac{1+\prod_{i=1}^{n}(1-2p_i)}{2}.
\]
You may use mathematical induction or repeated substitution in the recursion. State the initial value explicitly.

%\emph{Hint:} After obtaining the recursion, consider $2\Prob(A_n)-1$.
\item Under the independence assumption, is $\Prob(A_n)=1/2$ if and only if at least one $p_i=1/2$? Justify your answer. Then remove the independence assumption and specify probabilities for the four possible outcomes of two trials so that $p_1=p_2=1/2$ but $\Prob(A_2)\ne1/2$.
\end{enumerate}

\textbf{Solution: }

\begin{enumerate}
    \item[(a)]
    
    $A_2$ could be $SS$ and $FF$. $P(SS) = p_1 p_2$, $P(FF) = (1-p_1)(1-p_2)$, so

    $$
    P(A_2) = p_1 p_2 + (1-p_1)(1-p_2) = 1 - p_1 - p_2 + 2p_1 p_2 = \frac{1 + (1-2p_1)(1-2p_2)}{2}
    $$

    \item[(b)]
    
    $$
    P(A_n) = P(A_{n-1}) (1 - p_n) + (1 - P(A_{n-1})) p_n = P(A_{n-1}) (1 - 2p_n) + p_n
    $$

    $$
    P(A_n) - \frac{1}{2} = (1 - 2p_n) \left(P(A_{n-1}) - \frac{1}{2}\right)
    $$

    $$
    P(A_1) = 1 - p_1 = \frac{1 + (1 - 2p_1)}{2}
    $$

    $$
    P(A_k) - \frac{1}{2} = \frac{1}{2} \prod_{i=1}^k (1 - 2p_i)
    $$

    $$
    P(A_n) = \frac{1 + \prod_{i=1}^n (1 - 2p_i)}{2}
    $$

    \item[(c)]
    
    $$
    \frac{1 + \prod_{i=1}^n (1 - 2p_i)}{2} = \frac{1}{2} \iff \prod_{i=1}^n (1 - 2p_i) = 0
    $$

    So $\Prob(A_n) = 1/2$ if and only if at least one $p_i = 1/2$.

    Construct a counterexample for $n=2$ without independence.

    Let

    $$
    P(SS) = \frac{1}{2}
    $$

    $$
    P(SF) = P(FS) = 0
    $$

    $$
    P(FF) = \frac{1}{2}
    $$

    $$
    p_1 = P(SS) + P(SF) = \frac{1}{2}, \quad p_2 = P(SS) + P(FS) = \frac{1}{2}
    $$

    But $P(A_2) = P(SS) + P(FF) = \frac{1}{2} + \frac{1}{2} = 1 \neq \frac{1}{2}$.
\end{enumerate}
\end{homeworkProblem}

\newpage
\begin{homeworkProblem}[3]
% Source: Fall 2025 HW2, Problem 1(a)-(b), PDF page 2.
To battle against spam, Bob installs two anti-spam programs. An email is either legitimate (event $L$) or spam (event $L^c$). Let $M_j$ be the event that program $j$ marks the email as legitimate, for $j=1,2$. Assume that
\[
\Prob(L)=0.10,\qquad
\Prob(M_j\mid L)=\Prob(M_j^c\mid L^c)=0.90
\quad(j=1,2).
\]
The two programs process the same email. Their outputs are conditionally independent given $L$, and also conditionally independent given $L^c$.
\begin{enumerate}[(a)]
\item Find $\Prob(L\mid M_1)$ and simplify your answer.
\item Find $\Prob(L\mid M_1\cap M_2)$ and $\Prob(L\mid M_1\cap M_2^c)$. Compare each with the prior probability $\Prob(L)$.
\item Are $M_1$ and $M_2$ independent without conditioning on the email's true status? Prove your answer by computing the required probabilities. Explain why your answer is consistent with the stated conditional independence.
\end{enumerate}

\textbf{Solution: }

\begin{enumerate}
    \item[(a)]
    
    $$
    P(L \mid M_1)
    = \frac{P(M_1 \mid L) P(L)}{P(M_1 \mid L) P(L) + P(M_1 \mid L^c) P(L^c)}
    = 0.5
    $$

    \item[(b)]
    
    $$
    \begin{aligned}
    P(L \mid M_1 \cap M_2)
    &= \frac{P(M_1 \cap M_2 \mid L) P(L)}{P(M_1 \cap M_2 \mid L) P(L) + P(M_1 \cap M_2 \mid L^c) P(L^c)} \\
    &= \frac{P(M_1 \mid L) P(M_2 \mid L) P(L)}{P(M_1 \mid L) P(M_2 \mid L) P(L) + P(M_1 \mid L^c) P(M_2 \mid L^c) P(L^c)} \\
    &= 0.9
    \end{aligned}
    $$

    $$
    \begin{aligned}
    P(L \mid M_1 \cap M_2^c)
    &= \frac{P(M_1 \cap M_2^c \mid L) P(L)}{P(M_1 \cap M_2^c \mid L) P(L) + P(M_1 \cap M_2^c \mid L^c) P(L^c)} \\
    &= \frac{P(M_1 \mid L) P(M_2^c \mid L) P(L)}{P(M_1 \mid L) P(M_2^c \mid L) P(L) + P(M_1 \mid L^c) P(M_2^c \mid L^c) P(L^c)} \\
    &= 0.1
    \end{aligned}
    $$

    $P(L \mid M_1 \cap M_2) = 0.90 > P(L) = 0.10$, when both programs mark the email as legitimate, the probability that the email is legitimate increases.

    $P(L \mid M_1 \cap M_2^c) = 0.10 = P(L)$, when one program marks the email as legitimate and the other marks it as spam, the probability that the email is legitimate does not change.

    \item[(c)]
    
    According to $(a)$ and $(b)$, we have

    $$
    P(M_1) = 0.18, \quad P(M_2) = 0.18
    $$

    $$
    P(M_1 \cap M_2) = 0.09 \neg P(M_1) P(M_2)
    $$

    Therefore, $M_1$ and $M_2$ are not independent without conditioning.

    When observing $M_1$, the probability of $L$ increases from $0.1$ to $0.9$. And when the email is more likely to be legitimate, the probability of $M_2$ also increases. This is why $M_1$ and $M_2$ are not independent without conditioning.
\end{enumerate}
\end{homeworkProblem}

\newpage
\begin{homeworkProblem}[4]
% Mathematical source: Fall 2025 HW2, Problem 4(b)-(c), PDF page 5, with n=2.
% Extensions: optimal selective publication and the role of independence.
Consider a simplified model of an AI evaluation. In one run, an AI system answers two questions. Write $C$ for correct and $I$ for incorrect; the ordered outcomes $CC,CI,IC,II$ also denote the corresponding events. In parts (a) and (b), each answer is correct with probability $1/2$, independently of the other, so the four outcomes are equally likely.

An evaluator checks both answers against an answer key and may publish one of the following reports. Let $M_C$ and $M_I$ denote the events that the corresponding reports are published:
\[
\begin{array}{cl}
M_C:&\text{``At least one answer is correct'';}\\
M_I:&\text{``At least one answer is incorrect''.}
\end{array}
\]
At most one report is published per run, and every published report must be true.
\begin{enumerate}[(a)]
\item Initially, every run is reported. For $CC$, the evaluator publishes $M_C$; for $II$, the evaluator publishes $M_I$. For either $CI$ or $IC$, a fair coin is used to choose between the two reports, each with probability $1/2$. Using Bayes' rule, find
\[
\Prob(CC\mid M_C)\qquad\text{and}\qquad\Prob(II\mid M_I).
\]
\item Replace the rule in part (a) by any randomized reporting rule based on the outcome. The evaluator may choose not to publish a report for any of the four outcomes. The two reports must each have positive probability and must satisfy
\[
\Prob(CC\mid M_C)\ge\frac23,\qquad
\Prob(II\mid M_I)\ge\frac34.
\]
What is the smallest possible probability that no report is published? Specify a reporting rule for each of $CC,CI,IC,II$ that attains your answer, and prove that no admissible rule can do better.
\item Remove the independence assumption, but keep each answer's probability of being correct equal to $1/2$. The evaluator may choose a new randomized reporting rule, but must publish exactly one truthful report per run, with both report types having positive probability. Let $t=\Prob(CC)$. Determine all values of $t$ for which such a rule satisfies both posterior inequalities in part (b). Prove necessity and sufficiency, giving a valid rule for every claimed value of $t$. The rule may depend on $t$.
\end{enumerate}

\textbf{Solution: }

\begin{enumerate}
    \item[(a)]
    
    $$
    P(CC \mid M_C) = \frac{P(M_C \mid CC) P(CC)}{P(M_C \mid CC) P(CC) + P(M_C \mid CI) P(CI) + P(M_C \mid IC) P(IC)} = \frac{1}{2}
    $$

    $$
    P(II \mid M_I) = \frac{P(M_I \mid II) P(II)}{P(M_I \mid II) P(II) + P(M_I \mid CI) P(CI) + P(M_I \mid IC) P(IC)} = \frac{1}{2}
    $$

    \item[(b)]
    
    Obviously, $P(M_C \mid CC) =1$, $P(M_I \mid II) =1$.

    $$
    P(CC \mid M_C) = \frac{P(M_C \mid CC) P(CC)}{P(M_C \mid CC) P(CC) + P(M_C \mid CI) P(CI) + P(M_C \mid IC) P(IC)}
    $$

    $$
    P(II \mid M_I) = \frac{P(M_I \mid II) P(II)}{P(M_I \mid II) P(II) + P(M_I \mid CI) P(CI) + P(M_I \mid IC) P(IC)}
    $$

    $$
    \frac{1}{1 + P(M_C \mid CI) + P(M_C \mid IC)} \ge \frac{2}{3}
    $$

    $$
    \frac{1}{1 + P(M_I \mid CI) + P(M_I \mid IC)} \ge \frac{3}{4}
    $$

    Let $p_1$ be $P(M_C \mid CI)$, $p_2$ be $P(M_C \mid IC)$, $p_3$ be $P(M_I \mid CI)$, and $p_4$ be $P(M_I \mid IC)$. So the origin problem converts to:

    Subject to

    $$
    p_1 + p_2 \le \frac{1}{2}
    ,\quad
    p_3 + p_4 \le \frac{1}{3}
    ,\quad
    p_1 + p_3 \le 1
    ,\quad
    p_2 + p_4 \le 1
    $$

    maximize

    $$
    p_1+p_2+p_3+p_4
    $$

    So $\max{p_1+p_2+p_3+p_4} = \frac{5}{6}$, and $P(\text{no report}) = 1 - P(\text{report}) = 1 - \frac{1}{4}\big(P(M_C \mid CC) + P(M_I \mid II) + (p_1 + p_2 + p_3 + p_4)\big) = \frac{7}{24}$.

    One possible reporting rule is:

    $$
    P(M_C \mid CC) = 1
    ,\quad
    P(M_I \mid CC) = 0
    $$

    $$
    P(M_C \mid CI) = \frac{1}{4}
    ,\quad
    P(M_I \mid CI) = \frac{1}{6}
    $$

    $$
    P(M_C \mid IC) = \frac{1}{4}
    ,\quad
    P(M_I \mid IC) = \frac{1}{6}
    $$

    $$
    P(M_C \mid II) = 0
    ,\quad
    P(M_I \mid II) = 1
    $$

    \item[(c)]
    
    $$
    P(CI) = \frac{1}{2} - t, \quad P(IC) = \frac{1}{2} - t
    $$

    $$
    P(II) = t
    $$

    And $t \in \left[0, \frac{1}{2}\right]$.

    Let $P(M_C \mid CI) = p$, $P(M_C \mid IC) = q$.

    $$
    \begin{aligned}
    P(M_C)
    &= P(M_C \mid CC) P(CC) + P(M_C \mid CI) P(CI) + P(M_C \mid IC) P(IC) \\
    &= t + p \left(\frac{1}{2} - t\right) + q \left(\frac{1}{2} - t\right)
    \end{aligned}
    $$

    $$
    \begin{aligned}
    P(M_I)
    &= P(M_I \mid II) P(II) + P(M_I \mid CI) P(CI) + P(M_I \mid IC) P(IC) \\
    &= t + (1 - p) \left(\frac{1}{2} - t\right) + (1 - q) \left(\frac{1}{2} - t\right)
    \end{aligned}
    $$

    $$
    P(CC \mid M_C) = \frac{t}{t + (p + q) \left(\frac{1}{2} - t\right)} \ge \frac{2}{3}
    $$

    $$
    P(II \mid M_I) = \frac{t}{t + (2 - p - q) \left(\frac{1}{2} - t\right)} \ge \frac{3}{4}
    $$

    $t=0$ is impossible because $P(CC \mid M_C) = 0 < \frac{2}{3}$.

    $$
    \frac{t + (p + q) \left(\frac{1}{2} - t\right)}{t} \le \frac{3}{2}
    $$

    $$
    \frac{t + (2 - p - q) \left(\frac{1}{2} - t\right)}{t} \le \frac{4}{3}
    $$

    $$
    (p + q)\left(\frac{1}{2} - t\right) \le \frac{1}{2} t
    $$

    $$
    (2 - p - q)\left(\frac{1}{2} - t\right) \le \frac{1}{3} t
    $$

    $$
    2\left(\frac{1}{2} - t\right) \le \frac{5}{6} t
    $$

    $$
    t \ge \frac{6}{17}
    $$

    When $t = \frac{1}{2}$, we have $P(CI) = P(IC) = 0$, and

    $$
    P(CC \mid M_C) = 1 \ge \frac{2}{3}, \quad P(II \mid M_I) = 1 \ge \frac{3}{4}
    $$

    When $t \in \left[\frac{6}{17}, \frac{1}{2}\right)$, we have

    $$
    1 - \frac{t}{6\left(\frac{1}{2} - t\right)}
    \le \alpha
    \le \frac{t}{4\left(\frac{1}{2} - t\right)}
    $$

    $$
    \frac{t}{4\left(\frac{1}{2} - t\right)} - \left(1 - \frac{t}{6\left(\frac{1}{2} - t\right)}\right)
    = \frac{5t}{12\left(\frac{1}{2} - t\right)} - 1
    \ge \frac{5 \cdot \frac{6}{17}}{12 \cdot \frac{5}{34}} - 1
    = 0
    $$

    $$
    \frac{t}{4\left(\frac{1}{2} - t\right)} > 0
    , \quad
    1 - \frac{t}{6\left(\frac{1}{2} - t\right)} < 1
    $$

    So $\left[\max\left(0, 1 - \frac{t}{6\left(\frac{1}{2} - t\right)}\right), \min\left(1, \frac{t}{4\left(\frac{1}{2} - t\right)}\right)\right] \cap [0, 1]$ is non-empty.

    Let $\alpha = \min\left(1, \frac{t}{4\left(\frac{1}{2} - t\right)}\right)$.

    When $P(M_C \mid CI) = P(M_C \mid IC) = \alpha$,

    $$
    P(M_C) = t + 2\alpha\left(\frac{1}{2} - t\right) > 0
    , \quad
    P(M_I) = t + 2(1 - \alpha)\left(\frac{1}{2} - t\right) > 0
    $$

    $$
    P(CC \mid M_C)
    = \frac{t}{t + 2\alpha\left(\frac{1}{2} - t\right)}
    \ge \frac{t}{t + \frac{t}{2}}
    = \frac{2}{3}
    $$

    $$
    P(II \mid M_I)
    = \frac{t}{t + 2(1 - \alpha)\left(\frac{1}{2} - t\right)}
    \ge \frac{t}{t + \frac{t}{3}}
    = \frac{3}{4}
    $$

    Therefore,

    $$
    t \in \left[\frac{6}{17}, \frac{1}{2}\right]
    $$
\end{enumerate}
\end{homeworkProblem}

\end{document}
